\documentclass[slidetop,11pt]{beamer}
\usepackage[T1]{fontenc}
\usepackage[latin1]{inputenc}
\usepackage[frenchb]{babel}
% pour un pdf lisible � l'�cran si on ne dispose pas
% des fontes cmsuper ou lmodern
%\usepackage{lmodern}
\usepackage{aeguill}
%\hypersetup{pdfpagemode=FullScreen}

% ------------------------------------------------
%-----------   styles pour beamer   --------------
% ------------------------------------------------
\definecolor{fondtitre}{HTML}{336E17}
\definecolor{coultitre}{RGB}{105,13,13}
\xdefinecolor{fondtexte}{rgb}{1,0.95,0.86}     % ivoire
% -------------- Fioritures de style -------------
\beamertemplatetransparentcovered
% Mettre les icones de navigation en mode vertical (pour projection)
%\setbeamertemplate{navigation symbols}[vertical]
% ------------ Choix des th�mes ------------------
\usecolortheme{crane}
\setbeamercolor{bas}{fg=coultitre, bg=fondtitre!40}
\setbeamercolor{haut}{fg=fondtitre!40, bg=coultitre}
\beamerboxesdeclarecolorscheme{clair}{fondtitre!70}{coultitre!20}
\beamerboxesdeclarecolorscheme{compar}{coultitre!70}{fondtitre!20}
% ins�rer le nombre de pages
\logo{\insertframenumber/\inserttotalframenumber}
%------------ fin style beamer -------------------

% ----------- Contenu de la page de titre --------
\title{Grid in a Grid}
\subtitle{Deployment of a gLite Grid inside Grid'5000}
\author{Lucas Nussbaum - S�bastien Badia}
\institute{Grid'5000 - INRIA Grand-Est}
\date{g5ks \oldstylenums{Avril 2011}}
% ------------------------------------------------
\begin{document}
\frame{\titlepage}
%------------------ Sommaire ---------------
% 1 Intro
% 2 gLite Midelware
% 3 Demo
% 4 Next steps \LaTeX{}
\begin{frame}{Table of contents}
	\tableofcontents%[hideallsubsections]
\end{frame}

\section{Introduction}
\subsection{gLite}
\begin{frame}
\frametitle{Introduction}
\framesubtitle{gLite}
	\begin{block}{gLite (from wikip�dia)}
gLite is the middleware stack for grid computing used by the CERN LHC experiments and a very large variety of scientific domains. gLite is a part of the EGEE Project, gLite provides a complete set of services for building a production grid infrastructure.
	\end{block}
	\begin{center}
	  \includegraphics[width=4cm]{img/glite.png}
	\end{center}
\end{frame}

\subsection{Scientific Linux}
\begin{frame}
\frametitle{Introduction}
\framesubtitle{Scientific Linux}
	\begin{block}{Operating System}
		\begin{itemize}
			\item Scientific Linux 5.5 (Boron), based on RHEL
			\item Kadeploy3 use for deployment on Grid'5000
		\end{itemize}
	\end{block}
	\begin{center}
	  \includegraphics[width=8cm]{img/sl.jpg}
	\end{center}
\end{frame}

\subsection{Grid'5000}
\begin{frame}
\frametitle{Introduction}
\framesubtitle{Grid'5000}
	\begin{center}
	  \includegraphics[width=3cm]{img/g5k.png}
	\end{center}
	\begin{block}{Topologie}
		\begin{itemize}
			\item Multi-sites, Multi-Vo
			\item One Virtual Organisation per site.
		\end{itemize}
	\end{block}
\end{frame}

\subsection{gDeploy Script}
\begin{frame}
\frametitle{Introduction}
\framesubtitle{gDeploy Script}
	\begin{block}{Goal}
		\begin{itemize}
			\item Deploy a minimalist gLite site.
			\item composed by :
			\begin{itemize}
				\item a BDII element.
			  \item a Batch scheduler.
			  \item a Computing element.
			  \item Workers nodes.
			\end{itemize}
		\end{itemize}
	\end{block}
	\begin{block}{Script}
	  \begin{itemize}
			\item Written in ruby.
			\item avaliable on \url{http://sbadia.github.com/gdeploy/}
		\end{itemize}
	\end{block}
\end{frame}

\section{gLite Middleware}
\subsection{Information Service}
\begin{frame}
\frametitle{gLite Middleware}
\framesubtitle{Information Service}
  \begin{block}{gLite BDII}
	\begin{itemize}
		\item Information Service is a BDII (Berkley Database Information Index).
		\item BDII provide information about the grid ressources and their status.
	\end{itemize}
		\begin{beamerboxesrounded}[scheme=clair]{}
		  IS is a simple OpenLDAP server.
		\end{beamerboxesrounded}
	\end{block}
\end{frame}

\subsection{CE - Batch}
\begin{frame}
\frametitle{gLite Middleware}
\framesubtitle{CE - Batch}
  \begin{block}{Computing Element}
	\begin{itemize}
		\item Store information about workers nodes.
		\item Interface with cluster (wn).
	\end{itemize}
		\begin{beamerboxesrounded}[scheme=clair]{}
	    Cream computing element (torque client, mysql, tomcat).
		\end{beamerboxesrounded}
	\end{block}
  \begin{block}{gLite Batch}
	\begin{itemize}
		\item Batch scheduler.
		\item Queue manager.
   \end{itemize}
		\begin{beamerboxesrounded}[scheme=clair]{}
		 Batch is a Torque server and a Maui scheduler.
		\end{beamerboxesrounded}
	\end{block}
\end{frame}

\subsection{Workers Nodes}
\begin{frame}
\frametitle{gLite Middleware}
\framesubtitle{Workers Nodes}
  \begin{block}{Workers}
	\begin{itemize}
		\item Cluster on Scientific Linux 5.5
		\item belong to a Virtual Organisation.
	\end{itemize}
	\end{block}
	\begin{center}
	  \includegraphics[width=3cm]{img/wn.jpg}
	\end{center}
\end{frame}

\section{Demonstration}
\subsection{Deploy Scientific Linux}
\begin{frame}
\frametitle{Demonstration}
\framesubtitle{gDeploy}
	\begin{exampleblock}{Demo}
	  \begin{itemize}
	    \item Reserve and deploy Scientific Linux on nodes.
	    \item Launch gdeploy script.
	    \item Test your gLite site.
	  \end{itemize}
	\end{exampleblock}
\end{frame}

\section{Conclusion}
\subsection{Next Steps}
\begin{frame}
\frametitle{Conclusion}
\framesubtitle{Next steps}
Future work for gLite on Grid'5000.
	\begin{alertblock}{To do}
		\begin{itemize}
    	\item Generic sl image.
			\item SE element and Lfc.
			\item Inter Vo communications.
			 \item Wms and Ui.
		\end{itemize}
  \end{alertblock}
\end{frame}
\end{document}
